\section*{Hoofdstuk \ref{ch:b1:hotspots} --- Hotspots en intraplaatvulkanisme}

\begin{solution}{exo:b1:hotspots:1}
(a) $4205 + 5274 = \qty{9479}{m}$, bijna \qty{9.5}{km}. De vulkaan is
zwaarder dan het water dat zij verdringt: de plaat buigt onder haar last naar beneden, zodat
haar voet dieper ligt dan de bodem veraf, en haar ware hoogte groter is.
(b) $74000/0.1 = \num{740000}$ jaar bij de aanvoer van de Kilauea;
$74000/0.025 = \num{3.0e6}$ jaar bij de gemiddelde snelheid van de keten. (c)
$0.12\times19 = \qty{2.3}{km^3}$, ongeveer $\num{3e-5}$ van Mauna Loa: drie
duizendsten van een procent.
\end{solution}

\begin{solution}{exo:b1:hotspots:2}
(a) $0.92\times\pi/180 = \qty{0.0161}{rad.Myr^{-1}}$. (b) In radialen per
miljoen jaar maal kilometer is $\omega R\sin\Delta$ in kilometer per
miljoen jaar, dat is millimeter per jaar: $0.016057\times6371 = 102.3$, dus
$v = 102.3\sin\Delta$, wat \qty{51.1}{mm.yr^{-1}} geeft bij \ang{30},
\qty{88.6}{mm.yr^{-1}} bij \ang{60}, \qty{102.3}{mm.yr^{-1}} bij \ang{90} en
\qty{97.3}{mm.yr^{-1}} bij \ang{108.0}. (c) De snelheid is evenredig met de
afstand $R\sin\Delta$ tot de as, het grootst op de grote cirkel op \ang{90}
van de pool; voorbij die cirkel komt het punt weer dichter bij de as, aan de
kant van de tegenovergestelde pool.
\end{solution}

\begin{solution}{exo:b1:hotspots:3}
(a) Deccan: $\num{e6}/\num{5e5} = \qty{2}{km}$. Centrale Atlantische Oceaan:
$\num{2e6}/\num{7e6} = \qty{0.29}{km}$, ongeveer \qty{290}{m}. (b) Het gemiddelde voor
de Deccan ligt dicht bij haar doorsnede van \qty{2.5}{km}: de provincie is een dik,
doorlopend plateau. (c) Het Centrale Atlantische oppervlak is dat van het hele gebied,
op vier continenten, waarover haar overblijfselen zijn verspreid: lava's, dijken en
sills, nu door een oceaan gescheiden en sterk geërodeerd. Het gesteente ligt op sommige
plaatsen geconcentreerd en is op het grootste deel van het gebied afwezig, dus het gemiddelde
zegt niets over één stapel.
\end{solution}

\begin{solution}{exo:b1:hotspots:4}
(a) $\rho_m\alpha\,\Delta T = 3300\times\num{3.2e-5}\times250 =
\qty{26.4}{kg.m^{-3}}$. (b) De opwellingskolom heeft de extra massa
$(\rho_m - \rho_w)h = 2300\times690 = \qty{1.587e6}{kg.m^{-2}}$, die het lichte
laag in evenwicht moet houden: $H = \num{1.587e6}/26.4 = \qty{6.0e4}{m}$, ongeveer
\qty{60}{km} heet gesteente.
\end{solution}

\begin{solution}{exo:b1:hotspots:5}
(a) Bij A ($\lambda = 0$, $\varphi = 0$) is $\vec e = (0, 1, 0)$ en
$\vec n = (0, 0, 1)$; met $\sin\ang{270} = -1$ en $\cos\ang{270} = 0$ is
$\vec t_A = -\vec e = (0, -1, 0)$. Bij B ($\lambda = \ang{90}$) is
$\vec e = (-1, 0, 0)$, $\vec n = (0, 0, 1)$ en $\vec t_B = (1, 0, 0)$.
(b) $\vec t_A\times\vec t_B = (0, 0, 1)$. (c) Met $\vec p' = (0, 0, 1)$ en
$\vec r_A = (1, 0, 0)$, $\vec p'\times\vec r_A = (0, 1, 0)$, en
$(\vec p'\times\vec r_A)\cdot\vec t_A = -1 < 0$: de pool is
$-\vec p' = (0, 0, -1)$, de Zuidpool. Een westwaartse beweging langs de evenaar
is met de klok mee gezien vanuit boven de Noordpool, dus tegen de klok in gezien vanuit boven de Zuidpool.
(d) Op \ang{90} van de pool is $\omega = 50/6371 =
\qty{7.85e-3}{rad.Myr^{-1}} = \qty{0.45}{\degree\per\Myr}$.
\end{solution}

\begin{solution}{exo:b1:hotspots:6}
(a) $\sum x_i v_i = 90000 + 240000 + 342000 = \num{672000}$ en
$\sum x_i^2 = \num{7.0e7}$, dus $\omega = \qty{0.0096}{rad.Myr^{-1}}$, of
\qty{0.55}{\degree\per\Myr}. (b) De voorspellingen zijn 28.8, 48.0 en
\qty{57.6}{mm.yr^{-1}}; de residuen $+1.2$, $0.0$ en $-0.6$
\unit{mm.yr^{-1}}. (c) Voor een starre rotatie is de snelheid evenredig met de
afstand tot de as, en een punt op de as beweegt niet: er is geen snijpunt om te passen.
\end{solution}

\begin{solution}{exo:b1:hotspots:7}
(a) $v = \qty{50e-3}{m}/\qty{3.156e7}{s} = \qty{1.58e-9}{m.s^{-1}}$, dus
$B = 2300\times\num{3e8}\times\num{1.58e-9} = \qty{1090}{kg.s^{-1}}$ en
$Q = (910/\num{3.2e-5})\times1090 = \qty{3.1e10}{W}$. (b) Een derde van de
Hawaïaanse flux. (c) $B$ groeit met $A\,v$: de pluim onder de snellere plaat
moet dezelfde opwelling twee keer zo snel opnieuw maken, en haar flux is twee keer zo groot.
\end{solution}

\begin{solution}{exo:b1:hotspots:8}
(a) Dat is $350/(16.1 - 10.3) = 350/5.8 = \qty{60}{km.Myr^{-1}}$, dat is
\qty{60}{mm.yr^{-1}}. (b) Het jonge spoor, \qty{29 \pm 5}{mm.yr^{-1}}, stemt
overeen met de \qty{26 \pm 5}{mm.yr^{-1}} van de plaat binnen de onzekerheden; het
oude deel ging meer dan twee keer zo snel, $60 - 26 = 34$, bijna zeven
standaardonzekerheden van de plaatsnelheid. (c) Ja: N\,\ang{54}\,E is het \omterm{def:b1:spherical-astronomy:horizontal}{azimut}
\ang{54}, binnen één graad tegengesteld aan \ang{235}. Het vulkanisme beweegt over de
plaat de tegengestelde kant op ten opzichte van de plaat boven de \omterm{def:b1:hotspots:hotspot}{hotspot}. (d) De rek verplaatste de
oudere calderas uit elkaar, waardoor de afstanden tussen hen groter werden en zo de
schijnbare snelheid; de plaat zelf ging niet sneller.
\end{solution}

\begin{solution}{exo:b1:hotspots:9}
(a) $\num{9.2e10}/\num{e12} = \qty{0.092}{W.m^{-2}}$, \qty{92}{mW.m^{-2}}.
(b) Negen tot achttien keer de waargenomen overtollige warmte, die zelfs nul kan zijn.
(c) $L^2/\kappa =
\num{e10}/\num{e-6} = \qty{e16}{s}$, ongeveer \qty{320}{Myr}; de plaat steekt
\qty{1000}{km} over in $1000/92 = \qty{11}{Myr}$. Warmte van onderaf heeft tijd gehad om
slechts een klein deel van de plaat te doorkruisen terwijl de plaat boven de opwelling ligt: het grootste deel van de
warmte van de pluim blijft in de \omterm{def:b1:layered-earth:layers}{mantel} onder de plaat en wordt er mee weggevoerd, en het
oppervlak merkt er weinig van.
\end{solution}

\begin{solution}{exo:b1:hotspots:10}
(a) Tien keer de smelt, \qty{e7}{km^3}, met straal
$(3V/4\pi)^{1/3} = \qty{134}{km}$. (b) $\num{e7}/100 = \qty{e5}{km^2}$, een schijf van
$2\sqrt{\num{e5}/\pi} = \qty{357}{km}$ in diameter: achttien keer kleiner in oppervlak dan de
trappen. Het hoofd was groter of smolt minder, of het spreidde zich breder en dunner uit,
of de lava stroomde ver voorbij het hoofd. (c) \qty{1}{km^3} per jaar, tien
keer de huidige aanvoer van de Kilauea. (d) Een nieuwe pluim moet genoeg warm, licht gesteente
verzamelen om zich door de viskeuze \omterm{def:b1:layered-earth:layers}{mantel} omhoog te duwen, en een groot volume stijgt veel
sneller dan een klein; de staart die volgt stroomt omhoog langs het hete, zwakke pad dat het hoofd heeft geopend,
en een smalle buis volstaat.
\end{solution}

\begin{solution}{exo:b1:hotspots:11}
(a) $\qty{40}{km.Myr^{-1}}\times(81 - 47)\,\unit{Myr} = \qty{1360}{km}$.
(b) De plaat beweegt met $(-80\sin\ang{20}, 80\cos\ang{20}) = (-27.4, 75.2)$
en de \omterm{def:b1:hotspots:hotspot}{hotspot} met $(0, -40)$: ten opzichte van de \omterm{def:b1:hotspots:hotspot}{hotspot} beweegt de plaat met
$(-27.4, 115.2)$ \unit{mm.yr^{-1}}, een snelheid van \qty{118.4}{mm.yr^{-1}} naar
$\ang{360} - \arctan(27.4/115.2) = \ang{346.6}$. (c) Het spoor zou een
snelheid \qty{48}{\percent} te hoog en een richting \ang{6.6} te ver naar het oosten geven:
een wegdrijvende \omterm{def:b1:hotspots:hotspot}{hotspot} registreert de beweging van de plaat ten opzichte van zichzelf, niet ten opzichte van de diepe \omterm{def:b1:layered-earth:layers}{mantel}.
\end{solution}

\begin{solution}{exo:b1:hotspots:12}
(a) $\qty{3e6}{J.m^{-3}.K^{-1}}\times\qty{1300}{K}\times\qty{e15}{m^3} =
\qty{3.9e24}{J}$. (b) Verspreid over $\qty{3.16e13}{s}$ is dat \qty{1.2e11}{W}: ongeveer de
warmte die de Hawaïaanse pluim vandaag meevoert (1.3 keer), en \qty{0.26}{\percent}
van de \qty{47}{TW} van de Aarde. (c) Een ondergrens: het hoofd brengt ook de
overtollige warmte van de negen tiende van zijn gesteente dat niet smelt mee, en warmt de
basis van de \omterm{def:b1:layered-earth:lithosphere}{lithosfeer} op, en de latente warmte is buiten beschouwing gelaten.
\end{solution}

\begin{solution}{pb:b1:hotspots:1}
\textbf{1.} Op Hawaï ($\varphi = \ang{19.0}$, $\lambda = \ang{-155.2}$),
met $\sin\ang{304} = -0.8290$ en $\cos\ang{304} = 0.5592$:
\[
\begin{aligned}
  \vec r_H &= (-0.8583,\ -0.3966,\ 0.3256), &
  \vec e &= (0.4195,\ -0.9078,\ 0),\\
  \vec n &= (0.2955,\ 0.1366,\ 0.9455), &
  \vec t_H &= (-0.1825,\ 0.8289,\ 0.5287).
\end{aligned}
\]

\textbf{2.} Bij Louisville ($\varphi = \ang{-53.6}$, $\lambda = \ang{-140.6}$),
met $\sin\ang{316} = -0.6947$ en $\cos\ang{316} = 0.7193$:
\[
\begin{aligned}
  \vec r_L &= (-0.4586,\ -0.3767,\ -0.8049), &
  \vec e &= (0.6347,\ -0.7727,\ 0),\\
  \vec n &= (-0.6220,\ -0.5109,\ 0.5934), &
  \vec t_L &= (-0.8883,\ 0.1693,\ 0.4269).
\end{aligned}
\]

\textbf{3.} $\lVert\vec t_H\rVert^2 = 0.0333 + 0.6871 + 0.2795 = 0.9999$ en
$\vec t_H\cdot\vec r_H = 0.1566 - 0.3287 + 0.1721 = 0.0000$, binnen de afrondingsfout.

\textbf{4.} $\vec t_H\times\vec t_L = (0.2643, -0.3917, 0.7054)$, van lengte
$0.8491$. De hoek tussen de sporen is $\arcsin0.8491 = \ang{58.1}$ (hun inproduct,
$0.528$, is positief, dus de hoek is scherp).

\textbf{5.} $\vec p' = (0.3113, -0.4613, 0.8308)$. Dan is
$(\vec p'\times\vec r_H)\cdot\vec t_H = -0.98 < 0$: de pool is
$\vec p = (-0.3113, 0.4613, -0.8308)$.

\textbf{6.} $\arcsin(-0.8308) = \ang{-56.18}$ en
$\operatorname{atan2}(0.4613, -0.3113) = \ang{124.01}$: de pool ligt op
\ang{56.2}\,S, \ang{124.0}\,E, in de Zuidelijke Oceaan ten zuiden van Australië.

\textbf{7.} $\vec p\cdot\vec r_H = 0.2672 - 0.1830 - 0.2705 = -0.1863$, dus
$\Delta_H = \ang{100.7}$; $\vec p\cdot\vec r_L = 0.1428 - 0.1738 + 0.6687 =
0.6377$, dus $\Delta_L = \ang{50.4}$.

\textbf{8.} De vector $\vec p\times\vec r_L$ staat loodrecht op $\vec p$ en op $\vec r_L$, zoals $\vec t_L$ dat doet (door de constructie van $\vec p$, en omdat een spoor horizontaal is); beide liggen op dezelfde lijn, en de lengte ervan is $\sin\Delta_L$. Dus $(\vec p\times\vec r_L)\cdot\vec t_L = \pm\sin\Delta_L$, en de berekening geeft $+0.770 = \sin\ang{50.4}$: de plaat beweegt bij Louisville langs $+\vec t_L$, in de richting waarin het spoor loopt. De tekenregel, toegepast bij Hawaï, had bij Louisville de verkeerde zin kunnen geven; dat gebeurde niet, dus beide richtingen worden exact gereproduceerd, \ang{304} en \ang{316}.

\textbf{9.} $\omega = 67/(6371\sin\ang{50.38}) = 67/4908 =
\qty{0.01365}{rad.Myr^{-1}}$, dat is \qty{0.782}{\degree\per\Myr}.

\textbf{10.} $360/0.782 = \qty{460}{Myr}$.

\textbf{11.} De cosinusregel tussen \ang{56.18}\,S, \ang{124.01}\,E en
\ang{63.78}\,S, \ang{112.48}\,E geeft \ang{9.5}. De snelheid van de twee sporen,
\qty{0.78}{\degree\per\Myr}, ligt \qty{15}{\percent} onder de
\qty{0.92}{\degree\per\Myr} die op dertien snelheden is gepast.

\textbf{12.} Society \ang{73.6}, Pitcairn \ang{77.8}, Cobb \ang{134.7}.

\textbf{13.} $\omega R = 0.01365\times6371 = \qty{86.98}{mm.yr^{-1}}$, maal
$\sin\Delta$: Society \qty{83.5}{mm.yr^{-1}}, Pitcairn
\qty{85.0}{mm.yr^{-1}}, Cobb \qty{61.8}{mm.yr^{-1}}.

\textbf{14.} Pitcairn stemt overeen: \qty{5}{mm.yr^{-1}} minder dan gemeten, binnen
zijn \qty{15}{mm.yr^{-1}}. Society is \qty{25.5}{mm.yr^{-1}} te traag, 2.6
onzekerheden, en Cobb \qty{18.8}{mm.yr^{-1}} te snel, meer dan zes: geen van beide
stemt overeen.

\textbf{15.} De pool staat loodrecht op beide spoorvectoren. Als die bijna evenwijdig zijn, is hun kruisproduct kort, en een fout van enkele graden in een van beide richtingen draait het over een grote hoek: de pool is slecht bepaald.

\textbf{16.} $v = \qty{92e-3}{m}/\qty{3.156e7}{s} = \qty{2.92e-9}{m.s^{-1}}$
en $B = 2300\times\num{4.83e8}\times\num{2.92e-9} = \qty{3240}{kg.s^{-1}}$.

\textbf{17.} $Q = (910/\num{3.2e-5})\times3240 = \qty{9.2e10}{W}$,
\qty{0.20}{\percent} van \qty{47}{TW}.

\textbf{18.} De ruggen maken $21/67000 = \qty{3.1e-4}{km^3}$ per jaar per
kilometer; $0.1/\num{3.1e-4} = \qty{320}{km}$ rug.

\textbf{19.} $74000/0.1 = \num{740000}$ jaar.

\textbf{20.} $\num{2e6}/\num{81e6} = \qty{0.025}{km^3}$ per jaar: de Kilauea alleen
krijgt vandaag vier keer de gemiddelde productie van de hele \omterm{def:b1:hotspots:hotspot}{hotspot}, die bijna op haar sterkst is.

\textbf{21.} $v_H = 86.98\sin\ang{100.73} = 86.98\times0.9825 =
\qty{85.5}{mm.yr^{-1}}$.

\textbf{22.} $92 - 85.5 = \qty{6.5}{mm.yr^{-1}}$, 2.2 standaardonzekerheden van
de gemeten snelheid: de twee \omterm{def:b1:hotspots:hotspot}{hotspots} stemmen niet helemaal overeen over de beweging van één
\omterm{def:b1:plate-kinematics:rigid-plate}{starre plaat} over een vast referentiekader.

\textbf{23.} De leeftijden geven \qty{9.1}{cm} en \qty{8.6}{cm} per jaar, gemiddeld over de hele keten; de \qty{9.2}{cm} per jaar van Morgan en Morgan ligt dicht bij de eerste, en de voorspelling van Louisville, \qty{8.5}{cm} per jaar, dicht bij de tweede. Alle vier liggen binnen \qty{0.7}{cm} per jaar van elkaar: sporen, leeftijden en richtingen vertellen hetzelfde verhaal tot op ongeveer \qty{8}{\percent}.

\textbf{24.} Tot één decimaal, in centimeter per jaar, geeft het Louisville-spoor de Pacifische plaat boven de Hawaïaanse \omterm{def:b1:hotspots:hotspot}{hotspot} de snelheid
\textbf{\qty[text-series-to-math=true]{8.5}{cm.yr^{-1}}}.
\end{solution}
